% Presentación: Reconstrucción de MIP desde COU (UN Handbook F74 Rev.1, 2018)
% Compilar con pdflatex (Overleaf funciona directo).
\documentclass[aspectratio=169,10pt]{beamer}
\usepackage[utf8]{inputenc}
\usepackage[T1]{fontenc}
\usepackage[spanish,es-noshorthands]{babel}
\usepackage{amsmath,amssymb}
\usepackage{booktabs}
\usepackage{xcolor}
\usepackage{tikz}
\usetikzlibrary{arrows.meta,positioning}

\definecolor{cepal}{RGB}{31,78,121}
\definecolor{cepal2}{RGB}{46,117,182}
\definecolor{verdeok}{RGB}{112,173,71}
\definecolor{gris}{RGB}{89,89,89}

\usetheme{Madrid}
\setbeamercolor{structure}{fg=cepal}
\setbeamercolor{frametitle}{bg=cepal,fg=white}
\setbeamercolor{title}{fg=cepal}
\setbeamertemplate{navigation symbols}{}
\setbeamertemplate{itemize items}[circle]
\setbeamerfont{frametitle}{size=\large}

\newcommand{\hb}[1]{\textit{\color{gris}\small ``#1''}}
\newcommand{\dg}[1]{\widehat{#1}}

\title[MIP desde COU]{Reconstrucción de Matrices Insumo--Producto desde los COU}
\subtitle{Metodología paso a paso — UN Handbook on SUT and IOT (2018)}
\author{Equipo MIP}
\date{\today}

\begin{document}

\begin{frame}[plain]
  \titlepage
  \begin{center}\small\color{gris}
  Argentina 2022 como ejemplo — Naciones Unidas, \emph{Handbook on Supply, Use and Input--Output Tables}, Series F No.\,74 Rev.\,1
  \end{center}
\end{frame}

% ------------------------------------------------------------------
\begin{frame}{Qué hicimos, en una frase}
  \begin{block}{Objetivo}
  Convertir los \textbf{Cuadros de Oferta y Utilización (COU)} oficiales de INDEC en una
  \textbf{Matriz Insumo--Producto (MIP) simétrica}, siguiendo \emph{exactamente} el método del Handbook de la ONU.
  \end{block}
  \begin{itemize}
    \item Fuente única: los COU oficiales (precios corrientes).
    \item Cada identidad contable se verifica por construcción (error $\sim 10^{-15}$).
    \item Resultado auditable: los totales reconcilian celda a celda contra el COU.
  \end{itemize}
  \begin{alertblock}{Notación del Handbook que usaremos}
  $V$: oferta (industria$\times$producto) \quad $U$: utilización intermedia (producto$\times$industria) \quad
  $Y$: demanda final \quad $Z$: matriz simétrica de consumos intermedios
  \footnotesize\color{gris}\ (el Handbook la denota $S$ en producto$\times$producto y $B$ en industria$\times$industria).
  \end{alertblock}
\end{frame}

% ------------------------------------------------------------------
\begin{frame}{Punto de partida: el COU \small(dónde mirar en el archivo)}
  El COU de INDEC trae \textbf{dos hojas distintas} (así se llaman en el \texttt{.xls}):
  \begin{columns}
    \begin{column}{0.5\textwidth}
      \begin{block}{Oferta $V$ — hoja \texttt{Mat\_Of\_pc\_2022}}
      ¿Qué \textbf{produce} cada industria? A precios básicos. \\
      \footnotesize\color{gris} Filas = productos; columnas = industrias \emph{y además} columnas de valoración (importaciones, impuestos, márgenes, oferta total).
      \end{block}
    \end{column}
    \begin{column}{0.5\textwidth}
      \begin{block}{Utilización $U$ — hoja \texttt{Mat\_Ut\_pc\_2022}}
      ¿Qué \textbf{compra} cada industria? A precios de comprador. \\
      \footnotesize\color{gris} Filas = productos; columnas = industrias + demanda final.
      \end{block}
    \end{column}
  \end{columns}
  \vspace{0.4em}
  \begin{center}
  \textbf{Argentina 2022:}\quad $U$ es \textbf{222 productos $\times$ 107 industrias} (rectangular).
  \end{center}
  \begin{itemize}
    \item Filas = 222 \textbf{productos} (011 Cereales, 012 Semillas, animales vivos\ldots).
    \item Columnas = 107 \textbf{industrias} (011/014 Cultivos agrícolas, Cría de animales\ldots).
  \end{itemize}
\end{frame}

% ------------------------------------------------------------------
\begin{frame}{La meta: una MIP simétrica (cuadrada)}
  \begin{center}
  \begin{tikzpicture}[every node/.style={font=\small}]
    \draw[fill=cepal!10,draw=cepal] (0,0) rectangle (3.2,2.2);
    \node at (1.6,2.5) {\textbf{COU} $U$};
    \node at (1.6,1.1) {222 $\times$ 107};
    \node[below] at (1.6,0) {\color{gris}productos $\times$ industrias};

    \draw[-{Latex[length=3mm]},very thick,cepal] (3.6,1.1) -- (5.4,1.1)
      node[midway,above,align=center]{\footnotesize transformación\\\footnotesize (Cap.\,12)};

    \draw[fill=verdeok!15,draw=verdeok] (5.8,0) rectangle (8.0,2.2);
    \node at (6.9,2.5) {\textbf{MIP} $Z$};
    \node at (6.9,1.1) {107 $\times$ 107};
    \node[below] at (6.9,0) {\color{gris}industria $\times$ industria};
  \end{tikzpicture}
  \end{center}
  \begin{itemize}
    \item Una MIP es \textbf{cuadrada}: misma clasificación en filas y columnas.
    \item \textbf{¿Por qué cuadrada?} Para calcular la inversa de Leontief $L=(I-A)^{-1}$ y los multiplicadores hace falta una matriz cuadrada e invertible.
    \item No se puede invertir un rectángulo de $222\times107$.
  \end{itemize}
\end{frame}

% ------------------------------------------------------------------
\begin{frame}{¿Por qué la MIP no coincide celda a celda con el COU?}
  Porque responden a \textbf{preguntas distintas}:
  \begin{itemize}
    \item \textbf{Columna del COU:} ``¿cuánto compró la \emph{industria} Cría de animales?'' (todo junto).
    \item \textbf{Columna de la MIP:} ``¿cuánto le compró a la \emph{industria} Cultivos agrícolas?'' (por proveedor).
  \end{itemize}
  \begin{alertblock}{Clave}
  Son la \textbf{misma economía}, reorganizada de rectángulo ($222\times107$) a cuadrado ($107\times107$).
  \textbf{Nada se pierde}: los totales reconcilian exacto (lo veremos al final).
  \end{alertblock}
\end{frame}

% ------------------------------------------------------------------
\begin{frame}{El método del Handbook: 4 pasos (Modelo D)}
  \begin{enumerate}
    \item[\textbf{1}] \textbf{Participaciones de mercado} \hfill
      $\displaystyle D = V\,\dg{q}^{-1}$ \\
      {\footnotesize $D_{ip}$ = fracción del producto $p$ que produce la industria $i$ \; (sale de la Oferta).}
    \item[\textbf{2}] \textbf{Flujos intermedios} \hfill
      $\displaystyle Z = D\,U$ \quad {\footnotesize (industria $\times$ industria)}
    \item[\textbf{3}] \textbf{Coeficientes e inversa de Leontief} \hfill
      $\displaystyle A = Z\,\dg{g}^{-1}, \qquad L = (I-A)^{-1}$
    \item[\textbf{4}] \textbf{Demanda final doméstica} \hfill
      $\displaystyle f = D\,y^{\text{dom}}$
  \end{enumerate}
  \vspace{0.6em}
  \begin{block}{Lo que dice el Handbook sobre el Modelo D}
  \hb{Each product has its own specific sales structure, irrespective of the industry where it is produced. \textbf{No negative elements.}}
  \end{block}
\end{frame}

% ------------------------------------------------------------------
\begin{frame}{Ejemplo 1 — ¿Quién \emph{produce} cereales? (matriz $D$)}
  La matriz $D$ sale \textbf{solo de las columnas de industria} de la hoja de Oferta
  (\texttt{Mat\_Of\_pc\_2022}, fila \texttt{011 - Cereales}), a precios básicos:
  \begin{center}
  \begin{tabular}{lrr}
    \toprule
    Industria productora (columna del COU) & Producción & \% ($D$) \\
    \midrule
    011/014 \; Cultivos agrícolas & 2\,541.2 & \textbf{99.94\,\%} \\
    1531/2 \; Elaboración de prod.\ de molinería & 1.5 & 0.06\,\% \\
    1549 \; Tostado, molienda (café, etc.) & 0.03 & 0.00\,\% \\
    \midrule
    \textbf{Producción nacional de cereales} & \textbf{2\,542.7} & \textbf{100\,\%} \\
    \bottomrule
  \end{tabular}
  \end{center}
  \footnotesize\color{gris}\hfill (cifras en miles de millones de \$ corrientes)
  \begin{alertblock}{Corrección de precisión}
  \normalsize\color{black} Antes decíamos ``100\,\%''; el dato exacto del COU es \textbf{99.94\,\%}
  (más un 0.06\,\% de producción secundaria de la molinería, que \textbf{sí conservamos}). No es un supuesto: es lo que reporta el INDEC.
  \end{alertblock}
\end{frame}

% ------------------------------------------------------------------
\begin{frame}{Ejemplo 1 (cont.) — Producción $\neq$ Oferta total}
  \textbf{¿Por qué en el cuadro de oferta parece que ``otros sectores'' ofertan cereales?}
  Porque a la derecha de las industrias hay columnas que \textbf{no son productores} — son el puente de valoración:
  \begin{center}
  \begin{tabular}{lrr}
    \toprule
    Componente (columna de \texttt{Mat\_Of\_pc\_2022}) & Miles de mill.\ \$ & \% oferta \\
    \midrule
    Producción nacional a precios básicos (OPB) & 2\,542.7 & 92.4\,\% \\
    $+$ Impuestos a los productos netos (IP) & 185.0 & \phantom{0}6.7\,\% \\
    $+$ Márgenes de distribución (Mg) & 13.6 & \phantom{0}0.5\,\% \\
    $+$ Importaciones (CIF) & 7.2 & \phantom{0}0.3\,\% \\
    $+$ IVA no deducible \; $+$ derechos \; $+$ comisiones & 2.1 & \phantom{0}0.1\,\% \\
    \midrule
    \textbf{$=$ Oferta total a precios de comprador (OPC)} & \textbf{2\,750.6} & \textbf{100\,\%} \\
    \bottomrule
  \end{tabular}
  \end{center}
  \begin{itemize}\footnotesize
    \item Impuestos, IVA y derechos $\to$ gobierno; márgenes $\to$ comercio/transporte; importaciones $\to$ exterior. \textbf{Ninguno produce cereales}.
    \item La matriz $D$ usa \textbf{solo la fila OPB} (columnas de industria). El resto se trata aparte en la valoración (Cap.\,7).
  \end{itemize}
\end{frame}

% ------------------------------------------------------------------
\begin{frame}{Regla de agrupación 222 $\to$ 107 (Argentina 2022)}
  Cada compra de un \emph{producto} se atribuye a la(s) \emph{industria(s) que lo produce(n)}, según $D$.
  De los \textbf{221 productos con producción nacional}:
  \begin{center}
  \begin{tabular}{lr}
    \toprule
    Concentración del productor dominante & N.º de productos \\
    \midrule
    Un único productor (100\,\%) & 71 \\
    Productor dominante $\ge$ 99\,\% & 106 \\
    Productor dominante $\ge$ 95\,\% & 135 \\
    Productor dominante $\ge$ 90\,\% & 151 \\
    \midrule
    Con producción secundaria (2\,+ industrias) & 150 \\
    \bottomrule
  \end{tabular}
  \end{center}
  \begin{center}\small
  Participación dominante \textbf{media 89.6\,\%}, \textbf{mediana 98.4\,\%}: casi siempre hay un productor claro.
  \end{center}
\end{frame}

% ------------------------------------------------------------------
\begin{frame}{Ejemplo 2 — Maquila (producción secundaria repartida)}
  Otros productos los fabrican \textbf{muchas industrias}. Ej.: \emph{``Servicios de manufactura en insumos físicos''} (maquila),
  producido por \textbf{51 industrias}; uso intermedio total $= 496.6$ miles de mill.\ \$:
  \begin{center}
  \begin{tabular}{lrr}
    \toprule
    Industria que la produce & \% ($D$) & Se le atribuye \\
    \midrule
    Fabricación de otros productos & 14.3\,\% & 71.0 \\
    Hilatura, tejeduría y acabado  & 14.1\,\% & 70.0 \\
    Elaboración de aceites y grasas& 11.0\,\% & 54.6 \\
    Fabricación de prendas de vestir& \phantom{0}7.6\,\% & 37.7 \\
    Fabricación de vehículos automotores & \phantom{0}6.6\,\% & 32.8 \\
    \ldots (46 industrias más) & \ldots & \ldots \\
    \midrule
    \textbf{Suma} & \textbf{100\,\%} & \textbf{496.6} \\
    \bottomrule
  \end{tabular}
  \end{center}
  La compra se \textbf{reparte} entre las industrias productoras según $D$. \textbf{Nada se pierde} (los pedazos suman el total).
\end{frame}

% ------------------------------------------------------------------
\begin{frame}{El supuesto y por qué evita negativos}
  \begin{block}{Supuesto (para los productos repartidos)}
  Se asume que quien compró maquila la compró \textbf{proporcionalmente} a todas las industrias que la ofrecen (participaciones de mercado $D$). Es el método estándar del Handbook.
  \end{block}
  \begin{itemize}
    \item Para productos \textbf{concentrados} (cereales, 99.94\,\% de un productor) el supuesto es \textbf{irrelevante}: la atribución es prácticamente 1-a-1.
    \item Pesa solo en los \textbf{repartidos} (maquila y similares), porque el COU no registra a cuál productor le compró cada quien.
  \end{itemize}
  \begin{alertblock}{La gran ventaja: cero negativos}
  El reparto por participaciones de mercado (Modelo D) \textbf{nunca genera celdas negativas}.
  El método alternativo (tecnología de \emph{producto}, Modelos A/C) sí puede producir negativos
  \footnotesize\color{gris}(Cap.\,12, Anexo B) — es la causa de los ``valor agregado negativo'' que antes obligaban a excluir años.
  \end{alertblock}
\end{frame}

% ------------------------------------------------------------------
\begin{frame}{Valoración: de precios de comprador a básicos (Cap.\,7)}
  El $U$ del COU está a \textbf{precios de comprador} (con impuestos y márgenes) y mezcla doméstico + importado.
  La MIP debe estar a \textbf{precios básicos} y separar importaciones.
  \begin{enumerate}
    \item \textbf{Impuestos} sobre los productos $\to$ fila primaria aparte.
    \item \textbf{Márgenes} de comercio/transporte $\to$ se reasignan a los servicios que los proveen.
    \item \textbf{Importaciones} $\to$ fila primaria aparte (versión doméstica).
  \end{enumerate}
  \begin{block}{Lo que dice el Handbook}
  \hb{These differences must be eliminated either by proportional adjustments\ldots} (\S7.77) \\[2pt]
  Versión doméstica: \hb{national or domestic version\ldots IOT with endogenous imports\ldots the [imports] is inserted as a single row into the primary input part.} (Cap.\,12)
  \end{block}
\end{frame}

% ------------------------------------------------------------------
\begin{frame}{Reconciliación con el COU (la auditoría)}
  Aunque las celdas se reorganizan, \textbf{el total de compras de cada industria cuadra exacto} contra el COU.
  Ejemplo \textbf{Cultivos agrícolas (011/014), 2022}, miles de millones \$:
  \begin{center}
  \begin{tabular}{lr}
    \toprule
    COU: consumo intermedio de la industria (comprador) & \textbf{3\,340.2} \\
    \midrule
    $\Sigma\,Z$ (consumo intermedio doméstico, precios básicos) & 2\,753.8 \\
    $+$ importaciones intermedias & 411.0 \\
    $+$ impuestos y márgenes & 175.5 \\
    \midrule
    \textbf{= Suma} & \textbf{3\,340.2} \quad\color{verdeok}\textbf{(diferencia 0.00)} \\
    \bottomrule
  \end{tabular}
  \end{center}
  \vspace{0.3em}
  \begin{center}\small\color{gris}
  Es la hoja \textbf{``12. Auditoría COU''} del libro, con esta cuenta para las 107 industrias
  (columna a columna, diferencia $=0$).
  \end{center}
\end{frame}

% ------------------------------------------------------------------
\begin{frame}{Las identidades que se verifican}
  \begin{block}{Balance por fila (oferta = demanda)}
  \[ g_i = \sum_j z_{ij} + f_i \qquad \forall\, i \]
  \end{block}
  \begin{block}{Balance por columna (valor agregado)}
  \[ g_j = \sum_i z_{ij} + zm_j + W_j \qquad \forall\, j \]
  \footnotesize $zm_j$: consumo intermedio importado; \; $W_j$: valor agregado.
  \end{block}
  \begin{block}{Análisis}
  \[ A = Z\,\dg{g}^{-1}, \quad L = (I-A)^{-1}, \quad L\,f = g \]
  Validación: $a_{ij}\ge 0$, \; $a_{ij}\le 1$, \; $\sum_i a_{ij} < 1$ (hay valor agregado).
  \end{block}
  \begin{center}\small\color{verdeok}\textbf{Todas cierran con error $\sim 10^{-15}$ (precisión de máquina).}\end{center}
\end{frame}

% ------------------------------------------------------------------
\begin{frame}{Validación externa: Uruguay oficial}
  Uruguay \textbf{sí} publica su MIP oficial. La usamos como prueba de fuego:
  \begin{itemize}
    \item Su MIP oficial es \textbf{industria$\times$industria, precios básicos, importaciones en matriz aparte} $\Rightarrow$ \textbf{idéntica metodología a la nuestra}.
    \item Producción total: nuestro COU $=2{,}778{,}446$ vs.\ MIP oficial $=2{,}778{,}447$ \; \textbf{(100.0\,\% match)}.
    \item Su ``Efectos Directos e Inducidos'' es un modelo Tipo II (hogares endogenizados), no la inversa abierta.
  \end{itemize}
  \begin{block}{Por qué industria$\times$industria (Handbook, Cap.\,12)}
  \hb{Industry-by-industry IOTs are closer to statistical sources, business survey results and actual observations, and also to the SUTs.}
  \end{block}
\end{frame}

% ------------------------------------------------------------------
\begin{frame}{¿Qué es, exactamente, ``la MIP''?}
  \begin{center}
  \begin{tikzpicture}[every node/.style={font=\small,align=center}]
    \node[draw=cepal,fill=cepal!10,minimum width=3cm,minimum height=1.4cm] (z) {\textbf{Z}\\consumos\\intermedios};
    \node[right=0.3cm of z,draw=cepal,fill=cepal!10,minimum height=1.4cm] (f) {\textbf{f}\\dem.\\final};
    \node[below=0.2cm of z,draw=cepal,fill=cepal!10,minimum width=3cm] (w) {\textbf{W, zm} valor agregado, importado};
    \node[right=1.2cm of f,draw=verdeok,fill=verdeok!12,minimum width=3.2cm,minimum height=1.6cm] (d) {\textbf{A, L, B}\\multiplicadores};
    \draw[-{Latex}] (f.east) -- (d.west) node[midway,above]{\footnotesize se derivan};
  \end{tikzpicture}
  \end{center}
  \begin{itemize}
    \item \textbf{La MIP} propiamente dicha $=$ la matriz $Z$ con sus bordes ($f$, $W$, $zm$) que suman a $g$: la tabla contable.
    \item $A$, $L$, $B$, los multiplicadores $=$ objetos \textbf{derivados} para el análisis, no la MIP en sí.
  \end{itemize}
  \begin{center}\small\color{gris}
  La MIP es la \emph{fotografía} de la economía; $A$, $L$, $B$ son los \emph{lentes} para analizarla.
  \end{center}
\end{frame}

% ------------------------------------------------------------------
\begin{frame}{Resultados — Argentina (6 años)}
  \begin{center}
  \begin{tabular}{lccc}
    \toprule
    Año & Dimensión & fila $=$ columna & Multiplicador medio \\
    \midrule
    2004 & $162\times162$ & $7\times10^{-16}$ & 1.75 \\
    2018 & $107\times107$ & $1\times10^{-15}$ & 1.75 \\
    2019 & $107\times107$ & $7\times10^{-16}$ & 1.74 \\
    2020 & $107\times107$ & $8\times10^{-16}$ & 1.75 \\
    2021 & $107\times107$ & $1\times10^{-15}$ & 1.76 \\
    2022 & $107\times107$ & $1\times10^{-15}$ & 1.76 \\
    \bottomrule
  \end{tabular}
  \end{center}
  \begin{itemize}
    \item Cada año: un libro Excel con 13 pestañas (MIP, Z, vectores, diag($g$), balances, coeficientes $A$, validación, Leontief, $B$ y \textbf{Auditoría COU}).
    \item Mayores multiplicadores: agroindustria de cadena larga (carne, cuero, aceites, lácteos) — coherente con Argentina.
  \end{itemize}
\end{frame}

% ------------------------------------------------------------------
\begin{frame}{Resumen para llevarse}
  \begin{enumerate}
    \item El COU \textbf{basta} — es la fuente correcta y completa.
    \item La MIP es el COU \textbf{reorganizado} de rectángulo a cuadrado; \textbf{nada se pierde}.
    \item Agrupamos por \textbf{quién produce cada insumo} (matriz $D$, de la Oferta).
    \item Modelo D del Handbook: \textbf{sin negativos}, cercano a las fuentes, comparable con Chile/Colombia.
    \item Todo \textbf{reconcilia contra el COU} y cierra a precisión de máquina.
  \end{enumerate}
  \vfill
  \begin{center}
  \large\color{cepal}\textbf{Gracias — ¿preguntas?}
  \end{center}
\end{frame}

\end{document}
